% Preamble for transcriptions of Daniel Quillen's research notebooks.
%
% The fonds' preamble, whole — the same apparatus macros, the same tikz-cd set
% up, the same watermark — so that one day the same renderer can read both. What
% changes is the language (his notebooks are in English, and so is the edition)
% and the head of the document, which names the notebook and whose facsimile it
% was read from.
%
% A transcription sits one directory below transcripts-quillen/, as
% transcripts-quillen/<label>/batch-NN.en.tex, and opens with
%   % Preamble for transcriptions of Daniel Quillen's research notebooks.
%
% The fonds' preamble, whole — the same apparatus macros, the same tikz-cd set
% up, the same watermark — so that one day the same renderer can read both. What
% changes is the language (his notebooks are in English, and so is the edition)
% and the head of the document, which names the notebook and whose facsimile it
% was read from.
%
% A transcription sits one directory below transcripts-quillen/, as
% transcripts-quillen/<label>/batch-NN.en.tex, and opens with
%   % Preamble for transcriptions of Daniel Quillen's research notebooks.
%
% The fonds' preamble, whole — the same apparatus macros, the same tikz-cd set
% up, the same watermark — so that one day the same renderer can read both. What
% changes is the language (his notebooks are in English, and so is the edition)
% and the head of the document, which names the notebook and whose facsimile it
% was read from.
%
% A transcription sits one directory below transcripts-quillen/, as
% transcripts-quillen/<label>/batch-NN.en.tex, and opens with
%   % Preamble for transcriptions of Daniel Quillen's research notebooks.
%
% The fonds' preamble, whole — the same apparatus macros, the same tikz-cd set
% up, the same watermark — so that one day the same renderer can read both. What
% changes is the language (his notebooks are in English, and so is the edition)
% and the head of the document, which names the notebook and whose facsimile it
% was read from.
%
% A transcription sits one directory below transcripts-quillen/, as
% transcripts-quillen/<label>/batch-NN.en.tex, and opens with
%   \input{../preamble/quillen}
% The fonds' preamble reaches its tikz set-up by the same relative path, which
% is why tikz-setup.tex here only forwards to the real one.

\input{../../transcripts/preamble/grothendieck}

\newcommand{\notebook}[1]{\folder{#1}}
\renewcommand{\drawing}[1]{\par\smallskip{\small\color{gray!60!black}\textit{[Drawing: #1]}}\par\smallskip}

\renewcommand{\grothTitleBlock}{%
  \selectlanguage{english}%
  \begin{center}
    {\large\bfseries Daniel Quillen, research notebooks --- \grothFolder}\\[2pt]
    \ifx\grothFoldertitle\empty\else{\itshape\grothFoldertitle}\\[4pt]\fi
    {\small pages \grothFirst--\grothLast
      \ifx\grothBatch\empty\else\ (batch \grothBatch)\fi}\\[2pt]
    \ifx\grothDating\empty\else
      {\small\color{gray!60!black}Filed by the Clay under: \grothDating}\\[2pt]
    \fi
    \ifx\grothWatermark\empty\else
      {\small\bfseries\color{red!45!gray}\grothWatermark}\\[2pt]
    \fi
    {\footnotesize\color{gray!60!black}Transcribed from the facsimile published by the
     Clay Mathematics Institute.\\
     Uncertain readings are underlined, illegible ones marked [\,\ldots\,],
     additions in brackets.}
  \end{center}
  \vspace{1em}}

\endinput

% The fonds' preamble reaches its tikz set-up by the same relative path, which
% is why tikz-setup.tex here only forwards to the real one.

\input{../../transcripts/preamble/grothendieck}

\newcommand{\notebook}[1]{\folder{#1}}
\renewcommand{\drawing}[1]{\par\smallskip{\small\color{gray!60!black}\textit{[Drawing: #1]}}\par\smallskip}

\renewcommand{\grothTitleBlock}{%
  \selectlanguage{english}%
  \begin{center}
    {\large\bfseries Daniel Quillen, research notebooks --- \grothFolder}\\[2pt]
    \ifx\grothFoldertitle\empty\else{\itshape\grothFoldertitle}\\[4pt]\fi
    {\small pages \grothFirst--\grothLast
      \ifx\grothBatch\empty\else\ (batch \grothBatch)\fi}\\[2pt]
    \ifx\grothDating\empty\else
      {\small\color{gray!60!black}Filed by the Clay under: \grothDating}\\[2pt]
    \fi
    \ifx\grothWatermark\empty\else
      {\small\bfseries\color{red!45!gray}\grothWatermark}\\[2pt]
    \fi
    {\footnotesize\color{gray!60!black}Transcribed from the facsimile published by the
     Clay Mathematics Institute.\\
     Uncertain readings are underlined, illegible ones marked [\,\ldots\,],
     additions in brackets.}
  \end{center}
  \vspace{1em}}

\endinput

% The fonds' preamble reaches its tikz set-up by the same relative path, which
% is why tikz-setup.tex here only forwards to the real one.

\input{../../transcripts/preamble/grothendieck}

\newcommand{\notebook}[1]{\folder{#1}}
\renewcommand{\drawing}[1]{\par\smallskip{\small\color{gray!60!black}\textit{[Drawing: #1]}}\par\smallskip}

\renewcommand{\grothTitleBlock}{%
  \selectlanguage{english}%
  \begin{center}
    {\large\bfseries Daniel Quillen, research notebooks --- \grothFolder}\\[2pt]
    \ifx\grothFoldertitle\empty\else{\itshape\grothFoldertitle}\\[4pt]\fi
    {\small pages \grothFirst--\grothLast
      \ifx\grothBatch\empty\else\ (batch \grothBatch)\fi}\\[2pt]
    \ifx\grothDating\empty\else
      {\small\color{gray!60!black}Filed by the Clay under: \grothDating}\\[2pt]
    \fi
    \ifx\grothWatermark\empty\else
      {\small\bfseries\color{red!45!gray}\grothWatermark}\\[2pt]
    \fi
    {\footnotesize\color{gray!60!black}Transcribed from the facsimile published by the
     Clay Mathematics Institute.\\
     Uncertain readings are underlined, illegible ones marked [\,\ldots\,],
     additions in brackets.}
  \end{center}
  \vspace{1em}}

\endinput

% The fonds' preamble reaches its tikz set-up by the same relative path, which
% is why tikz-setup.tex here only forwards to the real one.

\input{../../transcripts/preamble/grothendieck}

\newcommand{\notebook}[1]{\folder{#1}}
\renewcommand{\drawing}[1]{\par\smallskip{\small\color{gray!60!black}\textit{[Drawing: #1]}}\par\smallskip}

\renewcommand{\grothTitleBlock}{%
  \selectlanguage{english}%
  \begin{center}
    {\large\bfseries Daniel Quillen, research notebooks --- \grothFolder}\\[2pt]
    \ifx\grothFoldertitle\empty\else{\itshape\grothFoldertitle}\\[4pt]\fi
    {\small pages \grothFirst--\grothLast
      \ifx\grothBatch\empty\else\ (batch \grothBatch)\fi}\\[2pt]
    \ifx\grothDating\empty\else
      {\small\color{gray!60!black}Filed by the Clay under: \grothDating}\\[2pt]
    \fi
    \ifx\grothWatermark\empty\else
      {\small\bfseries\color{red!45!gray}\grothWatermark}\\[2pt]
    \fi
    {\footnotesize\color{gray!60!black}Transcribed from the facsimile published by the
     Clay Mathematics Institute.\\
     Uncertain readings are underlined, illegible ones marked [\,\ldots\,],
     additions in brackets.}
  \end{center}
  \vspace{1em}}

\endinput
